\section*{Chapter \ref{ch:s1:momentum} --- Momentum}

\begin{solution}{exo:s1:momentum:1}
$0.12/0.40 = 0.3$; a raw loss of 34\% becomes $0.3 \times (-0.34) = -10.2\%$. That is April 2009 almost exactly: the weight was 0.29 and the scaled loss 10.1\%.
\end{solution}

\begin{solution}{exo:s1:momentum:2}
A gain of $1/0.5 - 1 = 100\%$. After a 50\% loss and a 100\% gain the strategy is back where it started: $0.5 \times 2 - 1 = 0$. Losses of momentum's size
take years of its average return to recover, which is why its maximum drawdown (78.4\%) matters more than its mean.
\end{solution}

\begin{solution}{exo:s1:momentum:3}
The weight traded is twice the one-way turnover: $2 \times 5.7 \times 0.001 = 1.14\%$ a year.
\end{solution}

\begin{solution}{exo:s1:momentum:4}
Over a decline, each stock's return is roughly its beta times the market's fall plus its own; the stocks that fell most, the past losers, are
disproportionately high-beta. Shorting them makes the book short beta. When the market rebounds sharply the high-beta losers rise most, and the book loses
on its shorts more than it gains on its lower-beta winners, at the moment when volatility is highest.
\end{solution}

\begin{solution}{exo:s1:momentum:5}
Over one month returns reverse (chapter~2). Including the last month would make the signal buy last month's winners, which tend to give some of it back,
diluting the momentum with a reversal bet of the opposite sign.
\end{solution}

\begin{solution}{exo:s1:momentum:6}
Scaling cuts exposure when recent volatility is high, which is when most crashes happen. It cannot see a crash that comes from a calm state: in September
1939 momentum's trailing volatility was low, the weight 0.92, and the crash arrived at almost full size. Scaling manages predictable risk, not all risk.
\end{solution}

\begin{solution}{exo:s1:momentum:7}
The residual book's Sharpe ratio goes from 0.67 to 0.70 over the scaled sample, the total-return book's stays at 0.52. On the French factor it goes from 0.45
to 0.87. Scaling helps in proportion to how much the strategy's volatility varies and how predictably: the French factor's volatility swings by a factor of
several between calm years and crises, while the synthetic books' varies much less.
\end{solution}

\begin{solution}{exo:s1:momentum:8}
Leverage multiplies the losses too: three times the raw factor's 2009 would have been a loss of more than 100\% of capital, and three times its maximum
drawdown of 78\% is ruin. A Sharpe ratio of 0.45 with a skewness of $-3$ cannot be levered like a normal return. Scale by volatility first, then choose
leverage on the scaled series, and keep the crash years in the sample.
\end{solution}

\begin{solution}{pb:s1:momentum:1}
\begin{enumerate}
\item Stocks with high returns over the past months keep outperforming; the 12--1 signal is the return from twelve months to one month ago.
\item Significant positive returns to buying winners and selling losers over 3- to 12-month holding periods, not explained by systematic risk, partly
reversed over the next two years.
\item 7.4\% a year, 16.3\% volatility, Sharpe ratio 0.45, skewness $-3.02$.
\item $-52.6\%$ in August 1932; $-78.4\%$ from June 1932 to September 1939.
\item Ranking on factor-model residuals; ranking industries on their returns. \omterm{def:s1:momentum:residual}{Residual momentum} earned about twice the risk-adjusted profits; \omterm{def:s1:momentum:residual}{industry
momentum} accounted for much of stock momentum and was highly profitable on its own.
\item ICs 0.020, 0.023 and 0.001; Sharpe ratios after costs 0.31, 0.45 and $-0.23$.
\item The same drift with less factor noise: a volatility of 11.0\% against 13.4\%.
\item The synthetic market plants no \omterm{def:s1:momentum:residual}{industry momentum}: its industry factors are independent from day to day.
\item Short periods of very large losses, in panic states after market declines and with high volatility, as the market rebounds.
\item After a decline the losers are high-beta stocks; the book is short beta, and a rebound hits it.
\item August and July 1932, April 2009, September 1939, January 2001, June 1938, June 1931, April 1933, November 2002 and January 2023.
\item Its market fell over the previous two years on only 6.7\% of days: too few bear markets.
\item \textbf{Named result.} In 2009 the French momentum factor lost 52.8\% (49.4\% from March to May); scaled to 12\% by its trailing 126-day volatility it
lost 16.8\% (16.5\%).
\item The Sharpe ratio rises from 0.45 to 0.87 and the skewness from $-3.02$ to $-0.40$.
\item When a crash comes from calm: September 1939, weight 0.92, a scaled loss of 29.1\%.
\item Leverage whenever momentum's volatility is below the target, with the financing and capacity it needs.
\item The market's first half-hour return predicts its last half-hour return, in S\&P 500 ETF data from 1993 to 2013 and in ten other ETFs.
\item The book's beta and its trend after market declines, the strategy's realised volatility, crowding measures, and turnover against costs.
\item Little: about 1.14\% a year against 8--38\% for the daily books.
\item Rare, very large losses when the market rebounds after a decline, which the average return does not show.
\end{enumerate}
\end{solution}

\begin{solution}{iq:s1:momentum:1}
Buying stocks that have outperformed over the past twelve months and selling those that have underperformed. The last month is skipped because at a
one-month horizon returns reverse.
\end{solution}

\begin{solution}{iq:s1:momentum:2}
After market declines, the losers the book is short are high-beta stocks; when the market rebounds sharply, often in high-volatility panic states, they rise
most and the book loses heavily: 1932, 2009.
\end{solution}

\begin{solution}{iq:s1:momentum:3}
Scale it by its forecast volatility; hedge or cap its market beta; reduce it after market declines with high volatility (Daniel and Moskowitz's panic state);
use \omterm{def:s1:momentum:residual}{residual momentum}; diversify with value and other strategies that do well in rebounds.
\end{solution}

\begin{solution}{iq:s1:momentum:4}
Total-return momentum carries last year's factor winners, so its factor exposures swing; \omterm{def:s1:momentum:residual}{residual momentum} ranks on the stock-specific part and is closer to
factor-neutral. \omterm{def:s1:momentum:residual}{Residual momentum}, for its steadier risk and, in the published evidence and here, a higher risk-adjusted return.
\end{solution}

\begin{solution}{iq:s1:momentum:5}
The market has fallen and the losers the book is short are now high-beta stocks: the book has become short the market. Hedge the beta or cut the book,
and watch for a rebound, which is when crashes happen.
\end{solution}

\begin{solution}{iq:s1:momentum:6}
With return $r_t = \mu + \sigma_t\epsilon_t$ and weight $w_t = c/\sigma_t$, the scaled return has mean $c\mu\,E[1/\sigma_t]$ and variance $c^2$ (plus a term
from $\mu^2$), so its Sharpe ratio is about $\mu E[1/\sigma_t]$, which by Jensen's inequality exceeds $\mu/\sqrt{E[\sigma_t^2]}$, the raw ratio. It lowers the
ratio when the mean is high exactly when volatility is high (scaling then cuts exposure when the strategy pays most) or when the forecast of volatility is poor.
\end{solution}
